% Corrected complete-order target and the Hirshfeld--Rabinovich obstruction.

\subsection{Why future until completeness stops at discrete orders}
\label{sec:p09-gpss-complete-order-complement}

GPSS Theorem~2.2 concerns discrete complete models. The modification
announced on its following page yields another route to Kamp's
\(S,U\) completeness result; it is not a theorem that \(U\)
alone expresses every future first-order formula on arbitrary complete
orders \cite[pp.~170--171]{GabbayPnueliShelahStavi1980}.
The latter statement is false even on the real line. We give the
Hirshfeld--Rabinovich nonexpressibility argument, including the ordinal
game used by its real-order construction
\cite[secs.~2--4 and appendix~A]{HirshfeldRabinovich2003Future}.
The source used here is the authors' version dated 6 March 2003;
its page locators refer to that sixteen-page version.
The source used here is the authors' version dated 6 March 2003;
its page locators refer to that sixteen-page version.

\begin{hbdefinition}{def:p09-hr-future-basis}{First-order future modalities and their depth}
A first-order truth table for an \(r\)-ary modality is a monadic
order formula \(\mu(t;X_1,\ldots,X_r)\) with only \(t\)
free as an individual variable. Its arguments are interpreted as
arbitrary unary sets. A truth table is \emph{future} when its value
at \(t\) agrees with its value in the induced labelled tail
\([t,+\infty)\), for every nonempty linear order and interpretation.
Equivalently it has a presentation with every quantifier restricted
to that tail. The forward implication follows by restricting all
quantifiers in the original formula; the reverse follows from the
semantics of that restriction.

The quantifier depth of a fixed truth table is the maximum number
of nested quantifiers. Write \(\mathsf{FTL}_d\), \(d\geq1\),
for the temporal language allowing every finite-arity future modality
with a supplied truth table of depth at most \(d\), together with
constants, proposition letters and Boolean connectives. A formula
has finite syntax but arbitrary modal nesting; the basis itself
may be infinite. The depth bound concerns individual truth tables,
not the first-order expansion of a whole temporal formula.
The strict until of \cref{def:p09-gpss-future-labels} belongs
to \(\mathsf{FTL}_2\).
\end{hbdefinition}

\begin{hbdefinition}{def:p09-hr-ordinal-models}{Finite powers of the natural order}
For \(k\geq1\), let \(O_k\) be the lexicographically ordered
set of \(k\)-tuples of natural numbers, with the leftmost
coordinate most significant. Its order type is \(\omega^k\).
Use finite Cantor polynomials
\[
                 \omega^{k-1}a_{k-1}+\cdots+\omega a_1+a_0
\]
for its points. The zero tuple is its least point.
Order sums and products in the following statements are ordinal
operations, not commutative numerical operations.

Two structures are \(n\)-equivalent, written \(\equiv_n\),
if they satisfy the same first-order sentences of quantifier
depth at most \(n\). An \(n\)-round comparison game chooses
one point in either structure at each round, answered by a
point in the other; the answering player wins when the
correspondence preserves equality, order and all unary
predicates on the selected points.
\end{hbdefinition}

\begin{hblemma}{thm:p09-hr-ordinal-game}{proved=annex:p9.hr.ordinal.game}{The complete ordinal comparison strategy}
Let \(n\geq0\), let \(\alpha,\beta\) be nonzero ordinals
below \(\omega^\omega\), and let \(\rho<\omega^n\).
The answering player has a winning \(n\)-round strategy on
\[
                   \omega^n\alpha+\rho,\qquad
                   \omega^n\beta+\rho.
\]
In particular \(O_k\equiv_n O_m\) whenever \(k,m\geq n\)
and \(k,m\geq1\). The annex gives the interval strategy
and proves its quantifier-depth consequence directly.
\end{hblemma}

\begin{hblemma}{thm:p09-hr-ordinal-tail}{proved=annex:p9.hr.ordinal.tail}{Future homogeneity of the ordinal models}
Every tail of \(O_k\) with an actual first point is isomorphic
to \(O_k\). Every future formula without unary predicates,
including a temporal formula made from constants and future
modalities, consequently has a constant truth value on \(O_k\).
For every \(\mathsf{FTL}_d\) formula without proposition
letters, this value is the same in \(O_k,O_m\) whenever
\(k,m>d+1\).
\end{hblemma}

\begin{hbdefinition}{def:p09-hr-future-limit-hierarchy}{A future hierarchy with its reference point retained}
For distinct current and tested variables \(x,t\), put
\[
\begin{aligned}
 L_0(x,t)&=(t=t),\\
 L_{j+1}(x,t)&=
 x<t\land
 \forall v\left(x\leq v<t\longrightarrow
       \exists u\,(v<u<t\land L_j(x,u))\right).
\end{aligned}
\]
All bound names are fresh. The fixed parameter \(x\) remains
the same throughout the recursion. Define
\[
                         \varphi_n(x)=\exists t\,(x<t\land L_n(x,t)),
                         \qquad n\geq1.
\]
These are future first-order formulas: every variable quantified
in a relevant branch is at or above \(x\). The explicit
condition \(x<t\) prevents a vacuous initial-point limit.

For a unary predicate \(P\), let \(\varphi_n^P(x)\) be
the formula obtained by restricting every quantified variable
in \(\varphi_n\) to \(P\), and conjoining \(P(x)\).
Existentials become \(\exists u(P(u)\land\cdots)\)
and universals become \(\forall u(P(u)\to\cdots)\);
all original order bounds remain. Thus
\(\varphi_n^P\) is still a future formula.
\end{hbdefinition}

\begin{hblemma}{thm:p09-hr-limit-separator}{proved=annex:p9.hr.limit.separator}{The exact threshold of the ordinal separators}
For \(k,n\geq1\),
\[
                         O_k,0\models\varphi_n
                                  \quad\Longleftrightarrow\quad k>n.
\]
More precisely, at a nonzero point of \(O_k\), \(L_n(0,t)\)
holds exactly when the last \(n\) Cantor coefficients of
that point are zero, with the convention that no nonzero
point below \(\omega^k\) has that property when \(n\geq k\).
For an ordered unary subset \(P\) and \(x\in P\),
\(\varphi_n^P(x)\) has exactly the value of \(\varphi_n\)
in the induced order on \(P\) at \(x\).
\end{hblemma}

\begin{hbdefinition}{def:p09-hr-continuous-models}{Real models with an ordinal set of marked points}
For \(k\geq1\), form the lexicographic order
\[
                C_k=\omega^k\times[0,1),\qquad
                P_k=\{(\gamma,0):\gamma<\omega^k\}.
\]
The only variable predicate is \(P_k\); every other proposition
letter, if used, is fixed false. The interval \([0,1)\)
here is real. These are ordered labelled structures, with
least point \((0,0)\).
\end{hbdefinition}

\begin{hblemma}{thm:p09-hr-continuous-models}{proved=annex:p9.hr.continuous.models}{Completeness, two tail types, and the lifted comparison game}
The order \(C_k\) is isomorphic to the nonnegative real line.
Its marked subset has type \(\omega^k\). Tails beginning at
marked points are mutually isomorphic, and tails beginning
at unmarked points are mutually isomorphic. Thus any
one-predicate future formula on \(C_k\) is equivalent
there to one of
\[
                              \top,\quad\bot,\quad P,\quad\neg P.
\]

If \(k,m\geq n\), \(k,m\geq1\), the labelled structures
\((C_k,P_k)\) and \((C_m,P_m)\) are \(n\)-equivalent.
The strategy plays the ordinal game in the first coordinate
and copies the chosen second coordinate exactly.
\end{hblemma}

\begin{hblemma}{thm:p09-hr-modal-invariance}{proved=annex:p9.hr.modal.invariance}{Arbitrary temporal nesting preserves the same two truth values}
For \(d\geq1\), every \(\mathsf{FTL}_d\) formula in the
single proposition letter \(P\) determines the same member
of \(\{\top,\bot,P,\neg P\}\) in every \(C_k\) with \(k>d\).
In particular no such formula distinguishes the least points
of \(C_k,C_m\) when \(k,m>d\).

The induction includes modalities of every finite arity and
arbitrarily deep nesting. In its modality step the operand
truth sets are replaced by their four possible unary
descriptions. The resulting truth table still has depth
at most \(d\); its uniform value on each of the two point
types is transferred by sentences of depth at most \(d+1\).
\end{hblemma}

\begin{hbtheorem}{thm:p09-hr-no-finite-future-basis}{proved=annex:p9.hr.no.finite.future.basis}{The finite-basis obstruction over real time}
For every \(d\geq1\), the future first-order formula
\(\varphi_{d+1}^{P}\) has no equivalent
\(\mathsf{FTL}_d\) formula over the real line.
It distinguishes \(C_{d+1}\) from \(C_{d+2}\) at their
least marked points, whereas the preceding lemma prevents
every such temporal distinction.

Consequently no finite set of first-order future modalities
is expressively complete for all future monadic first-order
formulas over the real line. The same conclusion holds on
the nonnegative reals, on the rational line and its
nonnegative part, and over the class of ordinals.
The finite-basis quantifier requires finite arities and
first-order future truth tables, as in the source
\cite[Theorems~3.6, 4.1 and 5.1]{HirshfeldRabinovich2003Future}.
\end{hbtheorem}

\begin{hbcorollary}{thm:p09-gpss-until-complete-order-obstruction}{proved=annex:p9.gpss.until.complete.order.obstruction}{An explicit counterinstance to until-only complete-order expressiveness}
The strict future-until language \(L(U)\) does not express
every future \(L_1\)-formula on complete orders.
The finite first-order formula \(\varphi_3^P\) is one
counterinstance: it is false at the least point of
\(C_3\) and true at the least point of \(C_4\),
while every \(L(U)\) formula in \(P\) agrees there.
The models may equivalently be placed on the full real
line and compared at zero.

Furthermore, even the enlarged left-labelled bounded
decomposition class of
\cref{def:p09-gpss-left-labelled-decomposition} cannot
be closed under all relative complements over complete
orders. If it were, its proved positive operations and
future-ray translation would give the false \(L(U)\)
completeness assertion. The obstruction therefore cannot
be removed by finishing a general final-gap recursion
while requiring all labels to remain in that class.
\end{hbcorollary}

\begin{hbexample}{def:p09-hr-until-table}{A three-quantifier future condition already separates until}
On every \(C_k\), the strict-until operation on the four truth sets
has this same table. Rows give the interval operand \(A\), columns
the witness operand \(B\), and entries give \(A\mathbin{\mathsf U_s}B\):
\[
\begin{array}{c|cccc}
 &\bot&\top&P&\neg P\\ \hline
 \bot&\bot&\bot&\bot&\bot\\
 P&\bot&\bot&\bot&\bot\\
 \top&\bot&\top&\top&\top\\
 \neg P&\bot&\top&\top&\top
\end{array}
\]
Every open nonempty interval has an unmarked point, and every point
has a next marked point with an entirely unmarked open gap before it.
These facts prove the table and are independent of \(k\).
Boolean induction therefore makes every until-only formula the
same one of the four truth sets on every \(C_k\).

Nevertheless the future formula
\[
\begin{split}
 P(x)\land\exists t\bigl(x<t\land P(t)\land
  \forall v\bigl((P(v)\land x\leq v<t)\to
        \exists u\,(P(u)\land v<u<t)\bigr)\bigr)
\end{split}
\]
is false at the least point of \(C_1\) and true at that of
\(C_2\). It is \(\varphi_1^P\): in the latter model the
marked point with index \(\omega\) is approached by the
marked points with finite indices. In \(C_1\), each
future marked point has a preceding marked point.
The appendix gives the table and this shorter counterinstance
within the proof of
\cref{thm:p09-gpss-until-complete-order-obstruction}.
\end{hbexample}

\begin{hbdefinition}{def:p09-hr-guarded-past-limit}{The valid guarded past-limit syntax}
The source's unary past-limit operation has truth table
\[
 K^-A(t)=
          \forall r\left(r<t\longrightarrow
                     \exists u\,(r<u<t\land A(u))\right).
\]
It is equivalent to \(\neg((\neg A)\mathbin{\mathsf S_s}\top)\),
where infix strict since puts its interval operand first.
This convention is vacuously true at a first point; it
is distinct from adding a nonemptiness conjunct.
\(K^-\) is a past-dependent operation.

Let \(\mathcal A\) contain all finite formulas built
using proposition letters, Boolean operations, strict
until and \(K^-\). Its \emph{guarded future fragment}
is the Boolean closure of proposition letters and
formulas \(B\mathbin{\mathsf U_s}C\), where
\(B,C\in\mathcal A\). Equivalently, every \(K^-\)
occurrence must lie under some strict until occurrence.
Thus its operands may contain a past-limit operation,
although every complete formula in the guarded
fragment is future. This is not the grammar \(L(U)\).
\end{hbdefinition}

\begin{hblemma}{thm:p09-hr-guarded-future}{proved=annex:p9.hr.guarded.future}{Why the guarded past-limit fragment is future}
Every formula of \(\mathcal A\) has the same truth
value at every \(t>x\) in a linear model and its
induced tail \([x,+\infty)\). Hence every formula
of the guarded future fragment is future at \(x\).
The strict guard is essential: \(K^-P\) itself
can distinguish two real-order interpretations
which agree at and after the current point.
\end{hblemma}

\begin{proofboundary}{def:p09-hr-positive-source-boundary}{The separate positive comparison}
The original paper's concluding discussion attributes
expressive completeness of this guarded \(U,K^-\)
fragment to Gabbay--Hodkinson--Reynolds, Theorem~10.3.20.
Its proof is not contained in the acquired paper.
The guarded-future lemma above proves its semantic scope.
The complete constructive equivalence over Dedekind-complete
orders is supplied in \cref{thm:p09-guarded-future-completeness},
using the distinct later primary presentation
\cite{Rabinovich2014Kamp}.
The original historical referral retains its own source identity
\cite[sec.~5, item~3]{HirshfeldRabinovich2003Future}.

GPSS's complete-order announcement remains the
\(S,U\) construction. The discrete \(U\)-only
theorems remain valid. The later almost-future
completeness result remains a distinct presentation.
Neither it nor the proved guarded comparison asserts
\(U\)-only completeness on arbitrary complete orders.
\end{proofboundary}
